\section{Theoretical Analysis}
\label{theoretical}

\paragraph{Motivation.}
Under the Dirichlet Compound Multinomial (DCM) model in \cite{choi2025debate}, the standard Bayesian debate induces a martingale over each agent's belief in the correct answer, suggesting that debate alone does not improve expected correctness.
A key gap is that real agents also apply \emph{private critique} after reading peers, and update their output based on their observation and critique~\cite{gou2023critic}.
We disentangle this effect with majority voting and show that
debate training improves multi-round debate by making this \emph{private critique} positively aligned with the ground truth,
thereby inducing a positive drift that breaks martingale neutrality.

\textbf{Setup.}
Fix an input question \(q\) and a finite answer set \(\mathcal{A}=\{1,\dots,K\}\). Without loss of generality, we set answer \(1\) as the correct one.
Following the DCM model~\cite{choi2025debate}, at round \(t\), agent \(i\) maintains Dirichlet parameters \(\alpha_{i,t}\in\mathbb{R}^K_{+}\) and observes neighbors \(N(i)\). Let $a(\cdot)$ extract the final answer choice from a response. We write $y_{i,t} := a(o_{i,t}) \in \mathcal{A}$ for the discrete answer label used in the DCM analysis.
We denote $L^1$-norm as \(\|v\|_1=\sum_{k=1}^K v^{(k)}\) and define the Dirichlet mean as \(\bar\theta_{i,t}:=\alpha_{i,t}/\|\alpha_{i,t}\|_1\in\Delta^{K}\).

\begin{definition}[Critique-augmented belief update]
\label{def:critique_update}
Let \(c_{i,t}\in\mathbb{N}^K\) be the neighbor count vector at round \(t\), 
\(c^{(k)}_{i,t}=\sum_{j\in N(i)}\mathbbm{1}\{y_{j,t-1}=k\}\).
In addition to neighborhood evidence \(c_{i,t}\), after observing the round-$t-1$ debate context, agent $i$ computes $\beta_{i,t-1}$
representing its \emph{private critique} from the debate context.
The belief update is
\begin{equation}
\alpha_{i,t} \;=\; \alpha_{i,t-1} \;+\; \beta_{i,t-1} \;+\; w_i\, c_{i,t},
\label{eq:update}
\end{equation}
where \(w_i\ge 0\) controls the strength of the social signal.
\end{definition}

\begin{definition}[Response generation via DCM]
\label{def:dcm}
Given \(\alpha_{i,t}\), the agent generates \(\theta_{i,t}\sim\mathrm{Dirichlet}(\alpha_{i,t})\)
and \(y_{i,t}\sim\mathrm{Categorical}(\theta_{i,t})\).
Marginally, \(\Pr(y_{i,t}=k\mid \alpha_{i,t})=\bar\theta^{(k)}_{i,t}\).
\end{definition}

Then the belief in the correct answer can be defined via Dirichlet mean:
\begin{equation}
p_{i,t} \;:=\; \bar\theta^{(1)}_{i,t} \;=\; \frac{\alpha^{(1)}_{i,t}}{\|\alpha_{i,t}\|_1}.
\label{eq:pt}
\end{equation}
We next let $\mathcal{F}_t$ denote the $\sigma$-algebra of information available at the start of round $t$ (see Appendix \ref{filtration} for the explicit definition).

\begin{definition}[Advantage of private critique]
\label{def:adv}
Assume \(\|\beta_{i,t}\|_1=m_\beta\) is constant (see Appendix \ref{beta_explained} for explanation).
The \emph{critique advantage} is
\begin{equation}
\delta_{i,t}\;:=\;\mathbb{E}\!\left[\beta^{(1)}_{i,t}\mid\mathcal{F}_{t}\right] \;-\; m_\beta\, p_{i,t}.
\label{eq:adv}
\end{equation}
Thus \(\delta_{i,t}>0\) means the critique allocates more belief to the correct answer in this round than a previous round's belief proportional to \(p_{i,t}\).
\end{definition}



\begin{theorem}[Critique induces belief drift]
\label{thm:drift}
Assume the mean-consistency condition \(\bar p_{N(i),t-1}=p_{i,t-1}\)
(the same condition under which standard MAD is a martingale in \cite{choi2025debate}). Let \(C_i:=m_\beta+w_i|N(i)|\), 
then
\begin{equation}
\mathbb{E}\!\left[p_{i,t}\mid \mathcal{F}_{t-1}\right]
=
p_{i,t-1}
+\frac{\delta_{i,t-1}}{\|\alpha_{i,t-1}\|_1+C_i}.
\label{eq:drift}
\end{equation}
\end{theorem}

\begin{proof}[Proof sketch]
Expand \(p_{i,t}\) from \eqref{eq:update} and take conditional expectation.
Use \(\mathbb{E}[c^{(1)}_{i,t}\mid\mathcal{F}_{t-1}]=\sum_{j\in N(i)}p_{j,t-1}=|N(i)|\bar p_{N(i),t-1}\)
and \(\alpha^{(1)}_{i,t-1}=p_{i,t-1}\|\alpha_{i,t-1}\|_1\) to obtain Lemma~\ref{lem:drift_decomp}.
Under mean-consistency assumption, the neighborhood drift cancels, yielding \eqref{eq:drift}.
\end{proof}

\begin{lemma}[Accumulated drift and diminishing returns]
\label{lem:accum_drift}
Assume mean-consistency and \(\delta_{i,t}\ge \mu\) for \(t=0,\dots,T-1\).
Let \(S_{i,0}:=\|\alpha_{i,0}\|_1\),
then
\begin{equation}
\begin{aligned}
\mathbb{E}[p_{i,T}]
&\ge p_{i,0}
+ \mu\sum_{t=1}^{T}\frac{1}{S_{i,0}+tC_i}
\ge p_{i,0}
+ \frac{\mu}{C_i}\log\!\left(\frac{S_{i,0}+(T+1)C_i}{S_{i,0}+C_i}\right).
\end{aligned}
\label{eq:accum_drift_main}
\end{equation}
\end{lemma}

\begin{proposition}[Training increases critique advantage]
\label{prop:training_adv}
Suppose debate training yields \(\delta_{i,t}\ge \mu>0\) on the evaluation-time debate distribution for early rounds.
Then Theorem~\ref{thm:drift} and Lemma~\ref{lem:accum_drift} imply \(\mathbb{E}[p_{i,t}]\) increases with \(t\).
\end{proposition}

\begin{lemma}[Correlation shrinks the effective ensemble size]
\label{lem:corr_neff}
Consider a single debate round with $N$ agents producing answers
$Y_1,\dots,Y_N \in \{1,\dots,K\}$.
Assume each $Y_n$ has the same marginal distribution $p\in\Delta^K$
with gap $\Delta := p_1-p_2>0$ (\cite{choi2025debate} Theorem 1).
Let $\hat p$ be the empirical distribution and
$y_{\mathrm{mv}}=\arg\max_k \hat p_k$ be the majority vote.

Let the correlation parameter be
\[
\rho := \max_{k\in[K]}\ \max_{a\neq b}\ 
\frac{\mathrm{Cov}(\mathbbm{1}\{Y_a=k\},\mathbbm{1}\{Y_b=k\})}{p_k(1-p_k)} \in [0,1].
\]
with the assumption that $p_k \in (0,1)$ for all $k$.
Then the plurality-vote error probability admits the bound
\[
\begin{aligned}
\Pr(y_{\mathrm{mv}}\neq 1)
&\le \frac{K(1+(N-1)\rho)}{N\,\Delta^2}
= \frac{K}{N_{\mathrm{eff}}\Delta^2},\\
N_{\mathrm{eff}} &:= \frac{N}{1+(N-1)\rho}.
\end{aligned}
\]
\end{lemma}

\paragraph{Role of majority vs.\ private critique.}
Our analysis disentangles the contributions of \emph{majority voting} and \emph{private critique} in multi-round debate.
When there is no training signal from critique (\emph{i.e.,} $\delta_{i,t}=0$), Theorem~\ref{thm:drift} reduces to a neutral (martingale) belief evolution, and the end-to-end accuracy is governed purely by vote amplification—recovering the standard majority-voting perspective (\emph{e.g.,} \cite{choi2025debate}).
\SDRL{} helps by increasing the critique advantage $\delta$, which induces positive drift in each agent’s correctness; however, Lemma~\ref{lem:accum_drift} shows that even under sustained advantage ($\delta_{i,t}\ge \mu$), the improvement accumulates only logarithmically in the number of rounds and yields diminishing per-round gains.
At the same time, as rounds progress agents condition on increasingly similar contexts, raising answer correlation and effectively reducing the ensemble size (Lemma~\ref{lem:corr_neff}), which weakens majority-vote amplification.
Together, these two effects explain a common empirical pattern: performance improves in early rounds (positive critique drift with low correlation) but peaks quickly and can decline later as marginal drift shrinks and homogeneity erodes the benefits of voting.

